\documentclass[12pt,a4paper]{article}

\usepackage[margin=2.5cm]{geometry}
\usepackage{amsmath,amssymb,amsthm}
\usepackage{booktabs,tabularx,array}
\usepackage{xcolor}
\usepackage{microtype}
\usepackage{enumitem}
\usepackage{listings}
\usepackage[numbers,sort&compress]{natbib}
\usepackage{hyperref}

\definecolor{codebg}{RGB}{247,248,250}
\definecolor{codeblue}{RGB}{28,82,135}
\definecolor{codegreen}{RGB}{38,115,75}
\lstdefinestyle{python}{
  language=Python,
  basicstyle=\ttfamily\small,
  keywordstyle=\color{codeblue}\bfseries,
  commentstyle=\color{codegreen},
  stringstyle=\color{purple!70!black},
  backgroundcolor=\color{codebg},
  frame=single,
  rulecolor=\color{black!18},
  breaklines=true,
  breakatwhitespace=false,
  columns=fullflexible,
  keepspaces=true,
  showstringspaces=false,
  tabsize=4,
  xleftmargin=0.5em,
  xrightmargin=0.5em
}

\hypersetup{
  colorlinks=true,
  linkcolor=blue!55!black,
  citecolor=green!45!black,
  urlcolor=blue!65!black,
  pdftitle={Optimal Mixing of Passive Scalars: Expanded Mathematical and Computational Guide},
  pdfauthor={Adebanji Adelowo}
}
\setlength{\emergencystretch}{3em}
\setlist[itemize]{leftmargin=*,itemsep=0.35em,topsep=0.45em}
\setlist[enumerate]{leftmargin=*,itemsep=0.45em,topsep=0.45em}

\newcommand{\norm}[1]{\left\lVert#1\right\rVert}
\newcommand{\Real}{\mathbb{R}}
\newcommand{\Torus}{\mathbb{T}^2}
\newcommand{\Lap}{\Delta}
\newcommand{\lapinv}{\Delta^{-1}}
\newcommand{\Proj}{P}
\newcommand{\dx}{\,dx\,dy}

\theoremstyle{definition}
\newtheorem{example}{Worked example}[section]
\newtheorem{remark}{Interpretation}[section]

\title{\textbf{Optimal Mixing of Passive Scalars}\\[5pt]
\large Expanded Mathematical, Algorithmic, and Computational Guide}
\author{Adebanji Adelowo}
\date{September 2026}

\begin{document}
\maketitle

\begin{abstract}
This companion report expands the mathematical and computational ideas in
the thesis \emph{Optimal Mixing of Passive Scalars}. It derives the
Lin--Thiffeault--Doering velocity step by step, explains the Fourier
pseudo-spectral discretisation from continuous operators to NumPy arrays,
and connects each equation to the corresponding Python implementation. It
also describes the adaptive Runge--Kutta solver, the conservation-based
stopping event, the extraction of mixing rates, and the interpretation of
the resolution study. The aim is not to replace the original thesis, but to
make its reasoning and code easier to study, reproduce, and extend.
\end{abstract}

\tableofcontents
\clearpage

\section{Problem in Plain Language}

A passive scalar is a quantity carried by a flow without changing that
flow. Dye in water is the standard example. The scalar field
$\theta(x,y,t)$ records the dye concentration, while $u(x,y,t)$ is the
velocity of the fluid. With no molecular diffusion, the governing equation
is
\begin{equation}
  \partial_t\theta+u\cdot\nabla\theta=0,
  \qquad \nabla\cdot u=0.
  \label{eq:transport}
\end{equation}
The first relation says that the concentration observed along a fluid
trajectory does not change. The second says that the flow is
incompressible: it rearranges area but does not create or destroy it.

Mixing therefore does not mean that $\theta$ becomes pointwise smaller.
Instead, stirring stretches large patches into increasingly thin positive
and negative filaments. A coarse observer sees those nearby values cancel,
even though the full distribution of scalar values is preserved.

The computational question is:
\begin{quote}
Given a fixed stirring budget, which incompressible velocity produces the
largest instantaneous decrease in a measure of unmixed large-scale
structure?
\end{quote}
The LTD strategy answers this question at every instant. It is a greedy
strategy: locally optimal in time, but not necessarily optimal over an
entire finite horizon \cite{LTD2011}.

\section{Why the \texorpdfstring{$H^{-1}$}{H-1} Norm Measures Mixing}

Let the periodic domain be $\Torus=[0,1]^2$, and write a mean-zero scalar as
\begin{equation}
  \theta(x)=\sum_{k\in\mathbb{Z}^2\setminus\{0\}}
  \widehat\theta_k e^{2\pi i k\cdot x}.
\end{equation}
Its squared $L^2$ norm is
\begin{equation}
  \norm{\theta}_{L^2}^2=\sum_{k\ne0}|\widehat\theta_k|^2.
\end{equation}
Every Fourier mode receives the same weight. In contrast,
\begin{equation}
  \norm{\theta}_{H^{-1}}^2
  =\sum_{k\ne0}\frac{|\widehat\theta_k|^2}{4\pi^2|k|^2}.
  \label{eq:hm1-fourier}
\end{equation}
The factor $|k|^{-2}$ strongly weights low wavenumbers and suppresses high
wavenumbers. Transferring scalar variance from a broad structure
($|k|$ small) to a fine filament ($|k|$ large) therefore reduces the
$H^{-1}$ norm without changing the $L^2$ norm.

\begin{example}
Suppose one unit of Fourier energy lies entirely in a mode with $|k|=1$.
Its contribution to \eqref{eq:hm1-fourier} is $1/(4\pi^2)$. Moving the same
energy to $|k|=10$ leaves the $L^2$ contribution equal to one, but reduces
the squared $H^{-1}$ contribution by a factor of $100$. This is precisely
the large-to-small-scale transfer that mixing is meant to detect.
\end{example}

If $\phi=\lapinv\theta$, so that $\Lap\phi=\theta$, integration by parts on
the torus gives an equivalent physical-space formula:
\begin{equation}
  \norm{\theta}_{H^{-1}}^2
  =-\int_{\Torus}\theta\phi\dx
  =\int_{\Torus}|\nabla\phi|^2\dx.
  \label{eq:hm1-physical}
\end{equation}
This identity is central to the velocity derivation.

\section{Deriving the Instantaneously Optimal Velocity}

\subsection{Step 1: Differentiate the objective}

Differentiate \eqref{eq:hm1-physical}. Since $\lapinv$ is self-adjoint,
the two differentiated terms are equal:
\begin{align}
  \frac{d}{dt}\norm{\theta}_{H^{-1}}^2
  &=-2\int_{\Torus}(\partial_t\theta)\lapinv\theta\dx \\
  &=2\int_{\Torus}(u\cdot\nabla\theta)\phi\dx,
\end{align}
where \eqref{eq:transport} was used in the second line. Incompressibility
and periodicity move the derivative from $\theta$ to $\phi$:
\begin{equation}
  \frac{d}{dt}\norm{\theta}_{H^{-1}}^2
  =-2\int_{\Torus}u\cdot
  \underbrace{(\theta\nabla\lapinv\theta)}_{g}\dx.
  \label{eq:first-variation}
\end{equation}

\subsection{Step 2: Keep only the divergence-free component}

The vector field $g$ need not be divergence-free. The Leray projection
\begin{equation}
  \Proj g=g-\nabla\lapinv(\nabla\cdot g)
  \label{eq:leray}
\end{equation}
removes its gradient component. For every divergence-free $u$,
$\int u\cdot\nabla q\dx=0$, so
\begin{equation}
  \int u\cdot g\dx=\int u\cdot\Proj g\dx.
\end{equation}
The discarded component cannot affect the objective under the admissible
velocity constraint.

\subsection{Step 3: Express the objective in the enstrophy inner product}

Define
\begin{equation}
  v=-\lapinv\Proj g.
  \label{eq:vdef}
\end{equation}
Then $\Proj g=-\Lap v$, and another integration by parts yields
\begin{equation}
  \frac{d}{dt}\norm{\theta}_{H^{-1}}^2
  =-2\int_{\Torus}\nabla u:\nabla v\dx.
  \label{eq:gradient-pairing}
\end{equation}
The colon denotes the Frobenius product: sum the products of all four
entries of the two velocity-gradient matrices.

\subsection{Step 4: Saturate the stirring budget}

The admissible velocities satisfy $\norm{\nabla u}_{L^2}\le F$. By
Cauchy--Schwarz,
\begin{equation}
  \int\nabla u:\nabla v\dx
  \le \norm{\nabla u}_{L^2}\norm{\nabla v}_{L^2}
  \le F\norm{\nabla v}_{L^2}.
\end{equation}
Equality holds when $u$ points in the same direction as $v$ in the
gradient inner product and uses the entire budget. Hence
\begin{equation}
  \boxed{u=F\frac{v}{\norm{\nabla v}_{L^2}},\qquad
  v=-\lapinv\Proj(\theta\nabla\lapinv\theta).}
  \label{eq:ltd}
\end{equation}
Substitution into \eqref{eq:gradient-pairing} gives
\begin{equation}
  \frac{d}{dt}\norm{\theta}_{H^{-1}}^2
  =-2F\norm{\nabla v}_{L^2}\le0.
\end{equation}
Thus the method cannot increase the mix norm unless numerical error
overwhelms the continuous argument.

\subsection{The degenerate case}

If $\Proj g=0$, then $v=0$ and the normalization in \eqref{eq:ltd} is
undefined. Mathematically, every admissible velocity has zero first-order
effect on the mix norm. The LTD paper treats this case through a
second-derivative optimization and an eigenvalue problem \cite{LTD2011}.
The implementation does not solve that secondary problem; it compares
$\norm{\nabla v}_{L^2}$ with $\norm{\lapinv g}_{L^2}$ and prints a warning
when their ratio becomes very small.

\section{From Fourier Analysis to NumPy Arrays}

\subsection{Grid and transform convention}

The code samples the unit torus on an $N\times N$ grid with spacing
$h=1/N$. NumPy's forward transform is unnormalized and its inverse includes
$1/N^2$. With the array indexing convention used in the code,
\begin{equation}
  \widehat f[k_y,k_x]
  =\sum_{i,j=0}^{N-1}f[i,j]
    e^{-2\pi i(ik_y+jk_x)/N}.
\end{equation}
Rows correspond to $y$ and columns to $x$. This small convention matters:
the $x$ multiplier must broadcast across columns, while the $y$ multiplier
must broadcast across rows.

\subsection{Wavenumbers and derivatives}

For even $N$, the FFT-ordered wavenumbers are
\begin{equation}
  0,1,\ldots,N/2-1,N/2,-N/2+1,\ldots,-1.
\end{equation}
The code constructs these values with
\begin{lstlisting}[style=python]
k = np.arange(N)
k_eff = (k - N * (k > N // 2)).astype(float)
kx = (2j * np.pi * k_eff)[np.newaxis, :]
ky = (2j * np.pi * k_eff)[:, np.newaxis]
\end{lstlisting}
For a Fourier mode $e^{2\pi i k\cdot x}$, differentiation multiplies its
coefficient by $2\pi i k_x$ or $2\pi i k_y$. Therefore an entire derivative
is computed by one elementwise multiplication in Fourier space.

For first derivatives, the code sets the Nyquist multiplier to zero. At
that frequency, the positive and negative representations coincide on an
even grid, so the sign required by an odd derivative is ambiguous. This is
a symmetry treatment, not a complete dealiasing method.

\subsection{Inverse Laplacian and zero mode}

Since
\begin{equation}
  \Lap e^{2\pi i k\cdot x}=-4\pi^2|k|^2e^{2\pi i k\cdot x},
\end{equation}
the inverse-Laplacian multiplier is
\begin{equation}
  \widehat{\lapinv f}_k=
  -\frac{\widehat f_k}{4\pi^2|k|^2},\qquad k\ne0.
\end{equation}
The zero mode is set to zero because constants lie in the kernel of the
Laplacian. This selects the mean-zero solution. In code:
\begin{lstlisting}[style=python]
lap = kx**2 + ky**2
LAP_INV = np.zeros_like(lap)
nz = lap != 0
LAP_INV[nz] = 1.0 / lap[nz]
\end{lstlisting}

\subsection{Discrete norm factors}

For grid samples on a unit square,
\begin{equation}
  \norm{f}_{L^p}\approx
  \left(h^2\sum_{i,j}|f_{ij}|^p\right)^{1/p}
  =h^{2/p}\norm{f}_{\ell^p}.
\end{equation}
Consequently the code multiplies the vector $\ell^2$, $\ell^4$, and
$\ell^8$ norms by $h$, $h^{1/2}$, and $h^{1/4}$, respectively. For the
unnormalized NumPy FFT, Parseval gives
\begin{equation}
  \norm{f}_{L^2}\approx\frac{\norm{\widehat f}_{\ell^2}}{N^2}.
  \label{eq:parseval-discrete}
\end{equation}
The factor $N^{-2}$ in the velocity normalization is therefore essential:
without it, the effective stirring budget would change with resolution.

\section{The Right-Hand Side Algorithm}

At each Runge--Kutta stage, the solver receives $\widehat\theta$ and must
return $\partial_t\widehat\theta$. Algorithmically, the calculation is:
\begin{enumerate}
  \item Compute $\phi=\lapinv\theta$ in Fourier space.
  \item Transform $\theta$, $\partial_x\phi$, and $\partial_y\phi$ to
        physical space.
  \item Multiply pointwise to obtain
        $g=(\theta\partial_x\phi,\theta\partial_y\phi)$, then transform
        both components back to Fourier space.
  \item Apply the Leray projection \eqref{eq:leray} mode by mode.
  \item Apply $-\lapinv$ to obtain $v$.
  \item Compute $\norm{\nabla v}_{L^2}$ using four spectral derivatives and
        the Parseval factor in \eqref{eq:parseval-discrete}.
  \item Normalize $v$ to obtain $u$ with
        $\norm{\nabla u}_{L^2}=F$.
  \item Transform $u$, $\partial_x\theta$, and $\partial_y\theta$ to
        physical space; form $u\cdot\nabla\theta$; transform back; negate.
\end{enumerate}

The core implementation follows the same order:
\begin{lstlisting}[style=python]
linv_th = LAP_INV * theta_hat
theta = np.real(np.fft.ifft2(theta_hat))
gx_hat = np.fft.fft2(
    theta * np.real(np.fft.ifft2(DEL_X * linv_th)))
gy_hat = np.fft.fft2(
    theta * np.real(np.fft.ifft2(DEL_Y * linv_th)))

div_g = DEL_X * gx_hat + DEL_Y * gy_hat
pgx = gx_hat - DEL_X * LAP_INV * div_g
pgy = gy_hat - DEL_Y * LAP_INV * div_g

vx_hat = -LAP_INV * pgx
vy_hat = -LAP_INV * pgy
\end{lstlisting}

The final transport step is
\begin{lstlisting}[style=python]
dtheta_hat = -np.fft.fft2(
    ux * np.real(np.fft.ifft2(DEL_X * theta_hat))
  + uy * np.real(np.fft.ifft2(DEL_Y * theta_hat)))
\end{lstlisting}

\subsection{Why products are formed in physical space}

Multiplication in physical space is convolution in Fourier space. A direct
two-dimensional convolution over all mode pairs costs $O(N^4)$. Two FFTs,
a pointwise multiplication, and an inverse FFT cost $O(N^2\log N)$, which
is why the method is called pseudo-spectral.

The speed comes with aliasing. Products create frequencies above the grid's
Nyquist limit, and those frequencies fold onto lower modes. The current code
does not apply explicit dealiasing. A standard improvement would retain the
lower two-thirds of the modes and use $3/2$ zero-padding for each quadratic
product. The conservation event detects accumulated resolution loss, but it
does not prevent aliasing before that threshold is crossed.

\section{State Representation and Adaptive Time Integration}

SciPy's \texttt{solve\_ivp} expects a real vector in the workflow used here,
whereas $\widehat\theta$ is complex. The code stores
\begin{equation}
  y=\begin{bmatrix}
     \operatorname{Re}(\widehat\theta)\\
     \operatorname{Im}(\widehat\theta)
  \end{bmatrix}\in\Real^{2N^2}.
\end{equation}
Each call to the right-hand side reconstructs the complex $N\times N$
array, performs the spectral calculation, and returns the real and imaginary
parts in the same flattened order.

The time integrator is RK45, the Dormand--Prince embedded pair. At each
candidate step it computes approximations of orders five and four. Their
difference estimates the local truncation error. The step is accepted when
the scaled error is below one and rejected otherwise; the next step size is
adjusted according to the observed error. Here
\begin{equation}
  \mathrm{rtol}=10^{-6},\qquad \mathrm{atol}=10^{-8}.
\end{equation}
The relative tolerance controls error proportional to the solution scale,
while the absolute tolerance protects components near zero.

\section{Resolution Check as an Event Function}

Exact incompressible transport preserves every $L^p$ norm. A numerical
change in these norms therefore signals spatial or temporal error. The code
monitors
\begin{equation}
  E(t)=\max\left(
  |L_2(t)-1|,
  \left|\frac{L_4(t)}{L_4(0)}-1\right|,
  \left|\frac{L_8(t)}{L_8(0)}-1\right|
  \right)-10^{-3}.
  \label{eq:event}
\end{equation}
SciPy searches for an upward zero crossing of $E$. When one is found, the
event is terminal and integration stops.

The higher norms serve a purpose. $L^8$ weights extreme localized values
more strongly than $L^2$, so it can reveal damage to sharp filaments sooner.
Using all three norms makes the diagnostic sensitive to different forms of
under-resolution.

\begin{remark}
The event is an error detector, not an error estimator with a rigorous bound
on every observable. A mixing rate fitted close to the stopping time can
still depend on resolution, which is exactly what the $N=32$ and $N=64$
comparison demonstrates.
\end{remark}

\section{Initial Data and Normalization}

Each initial-data function performs three conceptual operations:
\begin{enumerate}
  \item evaluate an analytic pattern on the grid;
  \item restrict or rearrange its support;
  \item divide by its discrete $L^2$ norm and shift it to the domain centre.
\end{enumerate}
For the sine bump,
\begin{equation}
  \theta_0(x,y)=
  \sin(2\pi x/a)\sin(2\pi y/a)
  \mathbf{1}_{[0,a)^2}(x,y),
\end{equation}
before normalization and centering. The parameter $a$ changes both the
support size and the characteristic spatial frequency. Smaller $a$ begins
with finer structure and reaches the grid scale sooner.

Normalization is implemented as
\begin{lstlisting}[style=python]
f /= np.linalg.norm(f.ravel()) * dx
\end{lstlisting}
because $dx=h$ and, in two dimensions,
$\norm{f}_{L^2}\approx h\norm{f}_{\ell^2}$.

\section{Extracting and Interpreting Mixing Rates}

If the normalized mix norm obeys
\begin{equation}
  M(t)=\frac{\norm{\theta(t)}_{H^{-1}}}
             {\norm{\theta_0}_{H^{-1}}}
  \approx e^{-rt},
\end{equation}
then
\begin{equation}
  \log M(t)\approx -rt.
\end{equation}
A least-squares line fitted to $\log M$ therefore has slope $-r$, and the
mixing timescale is $\tau=1/r$. The thesis fits the last two-thirds of each
trajectory to reduce the influence of the initial transient.

For $a=0.5$ and $N=64$, the reported late-window fit gives
$r\approx0.440$ and $\tau\approx2.27$. A full-trajectory fit gives
$r\approx0.394$. Both values are correct for their respective windows; the
difference shows that the curve is only approximately log-linear.

Across eight values of $a$, a second log-log regression fits
\begin{equation}
  r(a)\approx C a^{-\alpha}.
\end{equation}
Taking logarithms produces
\begin{equation}
  \log r\approx\log C-\alpha\log a,
\end{equation}
so $-\alpha$ is the slope of a linear fit to $(\log a,\log r)$. The study
finds $\alpha\approx1.61$ at $N=32$ and $\alpha\approx1.78$ at $N=64$.
Because the exponent changes materially with $N$, it is not yet numerically
converged. The two resolutions cannot establish its limiting value or
separate finite-resolution effects from the behavior of the greedy control.

\section{End-to-End Code Map}

\begin{table}[htbp]
\centering
\small
\caption{Relationship between mathematical tasks and Python functions.}
\begin{tabularx}{\linewidth}{@{}
  >{\raggedright\arraybackslash}p{3.2cm}
  >{\raggedright\arraybackslash}p{4.0cm}
  >{\raggedright\arraybackslash}X@{}}
\toprule
Task & Function & Main output \\
\midrule
Build grid and multipliers & \texttt{build\_operators} & Derivatives,
inverse Laplacian, mix-norm weight, and coordinate arrays \\
Construct $\theta_0$ & \texttt{idata\_sin}, \texttt{idata\_diag},
\texttt{idata\_strip}, \texttt{idata\_trigpoly} & Centered,
$L^2$-normalized physical array \\
Evaluate the ODE & \texttt{make\_convection\_hat} & Closure returning
$d\widehat\theta/dt$ in real-split form \\
Detect resolution loss & \texttt{make\_res\_check} & Terminal event based
on $L^2$, $L^4$, and $L^8$ drift \\
Integrate one case & \texttt{run\_simulation} & Times, fields, spectra,
and diagnostic norms \\
Fit decay rates & \texttt{replot\_norms} & Late-window slopes, timescales,
and scaling plots \\
\bottomrule
\end{tabularx}
\end{table}

The complete execution path is:
\begin{enumerate}
  \item call \texttt{build\_operators(N)} once;
  \item choose an initial-data function and a value of $a$;
  \item construct $\theta_0$, transform it, and pack the real state;
  \item build the right-hand side and resolution event as closures;
  \item call \texttt{solve\_ivp};
  \item unpack each saved state and compute diagnostic norms;
  \item repeat for the parameter sweep and fit the rates.
\end{enumerate}

\section{Numerical Risks and Recommended Extensions}

\begin{enumerate}
  \item \textbf{Add explicit dealiasing.} Implement $3/2$ padding for every
        quadratic product and compare against the present event-only
        safeguard.
  \item \textbf{Run a resolution ladder.} Use at least
        $N=32,64,128,256$ where feasible, and compare rates over a common
        physical fitting interval as well as each run's late window.
  \item \textbf{Quantify fit uncertainty.} Report standard errors or
        bootstrap intervals for $r$ and $\alpha$, and test sensitivity to
        the fitting window.
  \item \textbf{Guard the degenerate normalization.} Replace a diagnostic
        print with a documented fallback or controlled termination when
        $\norm{\nabla v}_{L^2}$ approaches machine precision.
  \item \textbf{Test invariants directly.} Add automated tests for the
        discrete divergence of $\Proj g$, the achieved stirring budget,
        Hermitian symmetry, mean preservation, and monotone initial
        decrease of the mix norm.
  \item \textbf{Compare control strategies.} Use finite-horizon or adjoint
        optimization to measure the performance gap between greedy and
        non-greedy stirring.
\end{enumerate}

\section{Summary}

The method can be understood as a short chain of ideas. The $H^{-1}$ norm
measures large-scale inhomogeneity. Differentiating it converts mixing into
an inner product between the velocity and
$\theta\nabla\lapinv\theta$. The Leray projection keeps the part that an
incompressible velocity can use. Applying $-\lapinv$ moves the problem into
the enstrophy inner product, and Cauchy--Schwarz identifies the steepest
admissible direction. FFT multipliers realize every differential operator,
while physical-space products evaluate the nonlinearity efficiently.
Adaptive RK45 advances the resulting ODE, and conservation drift determines
when the grid is no longer trustworthy.

This chain also clarifies the main limitation of the reported scaling law:
the fitted exponent changes with resolution because the usable trajectory
length changes with the grid. The present calculations establish a
reproducible framework and a clear finite-resolution result, but not an
asymptotic exponent.

\clearpage
\bibliographystyle{plainnat}
\begin{thebibliography}{9}
\bibitem{LTD2011}
Z.~Lin, J.-L.~Thiffeault, and C.~R.~Doering,
\newblock Optimal stirring strategies for passive scalar mixing,
\newblock \emph{Journal of Fluid Mechanics}, 675:465--476, 2011.

\bibitem{Iyer2014}
G.~Iyer, A.~Kiselev, and X.~Xu,
\newblock Lower bounds on the mix norm of passive scalars advected by
incompressible enstrophy-constrained flows,
\newblock \emph{Nonlinearity}, 27(5):973--985, 2014.

\bibitem{Trefethen2000}
L.~N.~Trefethen,
\newblock \emph{Spectral Methods in MATLAB},
\newblock SIAM, 2000.
\end{thebibliography}

\end{document}
